\section{The FATCD domain and the exact-six question}\label{sec:fatcd-domain}

% Mathematical job:
%   - Define finite audited typed closure descriptions (FATCD).
%   - Fix the neutral typed raw grammar (objects, states, relations, maps,
%     predicates, comparisons, paths, indices, orders, transitions,
%     refinements, traces, events, provenance, replay, check records, and
%     the Ann_* annotation kinds).
%   - State the finite-description, closure, and audit/bookkeeping conditions.
%   - State threshold attachment: annotations attach to raw features BEFORE
%     role projection is attempted.
%   - Define derived role projections pi_i(D) as interpretations of raw
%     features, NOT as domain membership slots.
%   - State the non-circularity and non-overbreadth theorem schemas.
%   - State the exact-six question precisely over FATCD.
%
% Notation to introduce (see paper/macros.tex once wired):
%   FATCD, D in FATCD, Raw(D), Audit(D), Closed(D),
%   Ann_L, Ann_Theta, Ann_Lift, Ann_Forget, Ann_Vis, Ann_Bridge,
%   Chan_i(D), Status_i(D), pi_i(D), Theta_i, W_i, Level_i,
%   Dec(D), Residual(D), Classify(r),
%   SixChan(Dec(D)), Exhaustive(Dec(D)), ScopedExactSix(FATCD).
%
% Governing sources:
%   workspace/meta/lens_lab/campaigns/paper3_exact_six_exhaustiveness/domain_exactification_api_repair/api_grounded_domain_repair_spec.md
%   workspace/meta/lens_lab/aristotle/results/paper3/paper3_exact_six_domain_exact_from_R008_B/paper3_exact_six_domain_exact_from_R008_B_result_shape.md
%   workspace/meta/lens_lab/papers/paper3/drafting_support/definition_statement_bank.md
%   workspace/meta/lens_lab/papers/paper3/drafting_support/notation_conventions.md
%   workspace/meta/lens_lab/papers/paper3/theorem_claims.md
%
% Overclaim risks:
%   - Writing the domain as if role membership is baked in.
%   - Treating optimization/objective data as excluded by definition.
%   - Collapsing "derived role projection" into "role occurrence".

The scoped exact-six theorem quantifies over a specific class of
admissible descriptions; this section defines that class. A
\emph{finite audited typed closure description}, abbreviated \FATCD{},
is a finitely presented collection of typed raw records and
annotations subject to finiteness, closure, and honest-bookkeeping
conditions, sitting in the broader finite-structure tradition of
finite model theory~\citep{Libkin2004}. Threshold annotations attach to raw features before any
role projection is attempted; the active projections
\(\pi_1(D),\dots,\pi_6(D)\) for the primitives \Pone{}--\Psix{} are
derived interpretations of raw material, not conditions of domain
membership. Candidate raw features not absorbed into such projections
are carried into an explicit residual classification, so seventh-role
threats are not excluded by definition. The section closes with
non-circularity and non-overbreadth schemas and formulates the
exact-six question precisely over \FATCD{}: whether every
\(D \in \FATCD{}\) admits a decomposition/exhaustion record with
exactly six role channels. The theorem chain answering that question
appears in
Sections~\ref{sec:lower-bounds}--\ref{sec:decomposition-exhaustion}.

The domain is framed in the finite theory-package setting of
Foundations~I~\citep{Tsiokos2026SixBirdsFoundations}, where one works
with tuples \(\mathcal{T} = (Z, f, \Sigma_f, E, \mathcal{A})\)
carrying a lens, a completion operator, and an audit functional. \FATCD{} does not
identify a description with such a tuple; it extracts from that
setting the structural requirements relevant to the scoped exact-six
question: finite typed raw data, closure under references, sources,
targets, endpoints, and annotations, and honest bookkeeping in the
sense of Section~\ref{sec:admissibility-meta-theory}. Optimization,
objective, and richer semantic material are not excluded at the
outset; they belong in the domain when finitely represented and
honestly recorded as raw or annotated data.

\subsection{Raw grammar}\label{subsec:raw-grammar}

We begin with the neutral typed raw grammar from which every element
of \FATCD{} is built. The alphabet is fixed in advance of the role
vocabulary and follows standard typed-record
conventions~\citep{CardelliMitchell1991}.

\begin{definition}[Raw record]\label{def:raw-record}
A \emph{raw record} is a finite typed syntactic or relational feature of a
description. The admissible record kinds are
\begin{itemize}[topsep=2pt,itemsep=1pt]
  \item \emph{object} and \emph{state} records;
  \item \emph{relation} records, \emph{map} records, and \emph{partial-map}
    records;
  \item \emph{predicate} records and \emph{comparison} records;
  \item \emph{path} or \emph{derivation} records;
  \item \emph{index}, \emph{order}, \emph{transition}, and \emph{refinement}
    records;
  \item \emph{trace}, \emph{event}, and \emph{provenance} records;
  \item \emph{replay} and \emph{check} records.
\end{itemize}
A raw record is not, by itself, a \Pone{}--\Psix{} role.
\end{definition}

\begin{definition}[Annotations]\label{def:annotations}
In addition to raw records, a description admits \emph{annotations} of six
declared kinds:
\begin{itemize}[topsep=2pt,itemsep=1pt]
  \item level annotations \(\mathrm{Ann}_L\);
  \item threshold annotations \(\mathrm{Ann}_\Theta\);
  \item lift annotations \(\mathrm{Ann}_{\mathrm{Lift}}\);
  \item forgetting annotations \(\mathrm{Ann}_{\mathrm{Forget}}\);
  \item instrument-visibility annotations \(\mathrm{Ann}_{\mathrm{Vis}}\);
  \item empirical-bridge annotations \(\mathrm{Ann}_{\mathrm{Bridge}}\).
\end{itemize}
These annotation kinds belong to the raw grammar; the finiteness and closure
conditions on how they attach to declared records are imposed in the next
subsection. For a description \(D\), we write \(\Raw{D}\) for the finite
collection of raw records and annotations of \(D\).
\end{definition}

\begin{remark}[Independence of the feature alphabet]\label{rem:alphabet-independence}
The record and annotation kinds in Definitions~\ref{def:raw-record} and
\ref{def:annotations} are syntactic or relational entities. No kind is named
\Pone{}, \ldots, \Psix{}. Consequently, membership in \FATCD{} does not
require any derived role projection, does not require an exact-six
decomposition/exhaustion record, and does not exclude candidate seventh-role
features by definition. The role vocabulary will enter only in
\S\ref{subsec:derived-role-projections} as a derived interpretation of raw
features under threshold and witness data.
\end{remark}

\subsection{Finiteness, closure, and audit conditions}\label{subsec:conditions}

Three further restrictions turn the raw grammar into a domain of
admissible descriptions: finiteness, closure, and audit.

\begin{definition}[Finite description]\label{def:finite}
A description \(D\) is \emph{finite} when
\begin{itemize}[topsep=2pt,itemsep=1pt]
  \item \(\Raw{D}\) contains finitely many records and annotations;
  \item every record or annotation carries finite typed fields;
  \item every relation, map, path, trace, or annotation is finitely
    represented;
  \item every reference in a record points to a declared finite record of
    \(D\);
  \item no infinite semantic, probabilistic, topological, geometric, causal,
    optimization, or resource structure is admitted unless it is finitely
    represented and bookkept as a raw feature of \(D\).
\end{itemize}
\end{definition}

\begin{definition}[Closed description]\label{def:closed}
A finite description \(D\) satisfies the \emph{closure condition} \(\Closed{D}\)
when
\begin{itemize}[topsep=2pt,itemsep=1pt]
  \item every source and target of every relation or map record is a declared
    record of \(D\);
  \item every endpoint of every path or derivation record is a declared record
    of \(D\);
  \item every index, transition, order, or refinement reference is a declared
    record of \(D\);
  \item every trace, event, provenance, replay, or check record references
    declared records of \(D\);
  \item every level, threshold, lift, forgetting, visibility, or
    empirical-bridge annotation references declared records of \(D\);
  \item every claim carried by \(D\) has a bookkeeping record.
\end{itemize}
\end{definition}

\begin{definition}[Audited description]\label{def:audited}
A description \(D\) is \emph{audited} when every claim carried by \(D\) has an
honest bookkeeping record. The collection of such records is denoted
\(\Audit{D} \subseteq \Raw{D}\), subject to the following discipline:
\begin{itemize}[topsep=2pt,itemsep=1pt]
  \item every promotion records its source level, target level, added data,
    losses, and nonclaims;
  \item every translation records its source context, target context,
    preserved data, suppressed data, and claim strength;
  \item every lift is marked as selected, not unique;
  \item every forgetting map records its losses;
  \item every visibility claim records both visible and suppressed data;
  \item every empirical-bridge claim records substrate, observable,
    instrument, level map, threshold witness relation, and
    suppressed-structure nonclaims.
\end{itemize}
\end{definition}

\subsection{Finite audited typed closure descriptions}\label{subsec:fatcd}

Combining the grammar with the three admissibility conditions yields
the domain over which the exact-six question is posed.

\begin{definition}[Finite audited typed closure description]\label{def:fatcd}
A \emph{finite audited typed closure description} is a finitely presented,
closed, audited collection of typed raw records and annotations, equipped with
honest bookkeeping for claims, levels, thresholds, visibility, forgetting,
lifts, and empirical-bridge assertions where present. Equivalently, it is a
description \(D\) whose raw content \(\Raw{D}\) is drawn from the kinds
declared in Definitions~\ref{def:raw-record} and~\ref{def:annotations}, such
that \(D\) is finite in the sense of Definition~\ref{def:finite}, satisfies
the closure condition \(\Closed{D}\) of Definition~\ref{def:closed}, and is
audited in the sense of Definition~\ref{def:audited}. We write \FATCD{} for
the class of all such descriptions, and \(D \in \FATCD{}\) for membership.
\end{definition}

\begin{remark}[No minimum richness]\label{rem:fatcd-minimality}
\FATCD{} imposes no lower bound on description size. A description
consisting of a single object record together with the bookkeeping
required by its claims may be admissible, as may a richer description
carrying transitions, paths or derivations, annotations of every kind,
and empirical-bridge annotations. The domain is defined by structural
discipline, not by the presence of any particular feature kind.
\end{remark}

\subsection{Threshold attachment}\label{subsec:threshold-attachment}

Membership in \FATCD{} is fixed; the next step is to specify how
role-theoretic claims may attach to the raw data of an admissible
description. Role projections in \FATCD{} enter only through a
disciplined attachment of threshold and witness data to raw features.
The condition below is the structural precondition for the role
vocabulary introduced in the next subsection.

\begin{definition}[Threshold attachment]\label{def:threshold-attachment}
A threshold annotation \(\mathrm{Ann}_\Theta \in \Raw{D}\) attaches first to
a finite cluster \(C \subseteq \Raw{D}\) of raw features of \(D\), where \(C\)
may consist of a single raw feature. No role projection is attempted on \(C\)
until all of the following are present in \(\Raw{D}\):
\begin{itemize}[topsep=2pt,itemsep=1pt]
  \item the raw features of \(C\) themselves;
  \item a threshold annotation \(\mathrm{Ann}_\Theta\) attached to \(C\);
  \item a witness schema specifying what in \(C\) would support the attempted
    projection;
  \item a record of the collapse cases that would invalidate the witness;
  \item a level annotation \(\mathrm{Ann}_L\) attached to \(C\);
  \item an audit record in \(\Audit{D}\) honestly bookkeeping the claim.
\end{itemize}
\end{definition}

\begin{remark}[Pre-projection discipline]\label{rem:pre-projection}
The items of Definition~\ref{def:threshold-attachment} form a structural
precondition, not a process narrative. What matters is that each item is
present explicitly in \(\Raw{D}\) before a role projection is asserted, so
that no threshold, witness, collapse case, or level is supplied tacitly by
the act of projection itself.
\end{remark}

\subsection{Derived role projections}\label{subsec:derived-role-projections}

The role vocabulary \Pone{}--\Psix{} now enters \FATCD{}, not as a
constraint on membership but as a derived interpretation of raw
features conditioned on the attachment data of
\S\ref{subsec:threshold-attachment}.

\begin{definition}[Derived role projection]\label{def:role-projection}
A \emph{derived role projection} for a description \(D \in \FATCD{}\) at
role index \(i \in \{1,\ldots,6\}\) is an active derived interpretation
\(\pi_i(D)\) of raw features of \(D\) as the role \(\mathrm{P}_i\), supported
by threshold data, witness schema, level assignment, and honest bookkeeping.
Concretely, \(\pi_i(D)\) is specified by
\begin{itemize}[topsep=2pt,itemsep=1pt]
  \item the proposed role name, one of \Pone{}, \Ptwo{}, \Pthree{},
    \Pfour{}, \Pfive{}, \Psix{};
  \item a raw feature cluster \(C \subseteq \Raw{D}\) on which the role is
    interpreted;
  \item the threshold data \(\Theta_i\) attached to \(C\) in the sense of
    Definition~\ref{def:threshold-attachment};
  \item a witness schema \(W_i\) specifying what constitutes a witness for
    the candidate role under its fixed meaning;
  \item a level assignment \(\mathrm{Level}_i\);
  \item an evidence or certificate type;
  \item a record of collapse cases avoided or accepted;
  \item loss and lift annotations for any level crossings involved;
  \item the explicit nonclaims preserved by the projection.
\end{itemize}
The supporting data for \(\pi_i(D)\) are carried by \(\Raw{D}\) and certified
by \(\Audit{D}\) in the sense of Definition~\ref{def:audited}.
\end{definition}

\begin{remark}[Fixed role meanings]\label{rem:fixed-role-meanings}
The role names \Pone{}--\Psix{} are defined by their meanings in
Section~\ref{sec:primitive-role-architecture}: \Pone{} concerns
update-vs-quotient compatibility and descent obstruction; \Ptwo{} concerns
macro-specification representability; \Pthree{} concerns enriched
route and coherence mismatch, not a directionality certificate; \Pfour{}
concerns stage, order, transition, and refinement structure; \Pfive{}
concerns quotient packaging; and
\Psix{} concerns audit, bookkeeping, provenance, replay, and checkability.
The witness schema \(W_i\) in Definition~\ref{def:role-projection} must be
aligned with these meanings; a cluster whose raw features do not support the
meaning of the proposed role contributes no \(\pi_i(D)\).
\end{remark}

\begin{remark}[Projection and membership are independent]\label{rem:projection-vs-membership}
Whether a description \(D\) supports a given derived role projection
\(\pi_i(D)\) is not part of the admissibility conditions of
Definitions~\ref{def:finite}--\ref{def:audited}. A description
\(D \in \FATCD{}\) may carry no derived role projection, projections
for some roles, or projections for all six; its membership in
\FATCD{} is unaffected. The decomposition/exhaustion record of
Section~\ref{sec:decomposition-exhaustion} accordingly assigns a
status to every role channel of \(D\), only some of which need be
active.
\end{remark}

\subsection{Residual classification scheme}\label{subsec:residual-classification}

Some raw features remain role-theoretically unsettled at the
raw-grammar level. \FATCD{} provides an explicit scheme for
classifying them.

\begin{definition}[Residual feature; residual classification scheme]\label{def:residual-scheme}
Let \(D \in \FATCD{}\). A \emph{residual feature} of \(D\) is a raw record
or annotation in \(\Raw{D}\) whose role-theoretic disposition is not fixed
solely by its raw-grammar kind. The \emph{residual classification scheme}
assigns to each residual feature \(r\) a single label
\(\operatorname{Classify}(r)\) drawn from the list
\begin{enumerate}[label=(\roman*),topsep=2pt,itemsep=1pt]
  \item already absorbed into a derived \Pone{}--\Psix{} projection
    \(\pi_i(D)\);
  \item composite of such projections;
  \item presentation variant of an already classified feature;
  \item level shift of an already classified feature;
  \item activation-threshold refinement;
  \item instrument-visibility artifact;
  \item empirical-bridge annotation;
  \item outside-domain structure;
  \item already-covered bookkeeping or annotation data;
  \item unresolved candidate seventh-role threat;
  \item genuine seventh irreducible typed role.
\end{enumerate}
The scheme is declarative. At the domain-definition level, \FATCD{} records
the available labels but neither proves that every residual feature lands in
categories (i)--(ix) nor proves that categories (x) and (xi) are empty. Those
claims are deferred to Section~\ref{sec:upper-bound-classification}.
\end{definition}

\begin{remark}[Optimization and objective data are not excluded]\label{rem:optimization-not-excluded}
Raw features expressing optimization, objective, or selection criteria
pass through the scheme of Definition~\ref{def:residual-scheme}; they
are not dismissed by definition. Under the fixed role meanings
(Remark~\ref{rem:fixed-role-meanings}), typical classifications
include: a predicate or specification record possibly supporting a
\Ptwo{} projection; a map or transition-selection record possibly
supporting a \Pone{} or \Pfour{} projection; a route-comparison or
route-choice record possibly supporting a \Pthree{} projection; an
audit or check record possibly supporting a \Psix{} projection; an
empirical-bridge annotation; or an outside-domain structure when the
feature requires non-finite or unbookkept semantics. A feature that
fits none of these is recorded under category~(x).
\end{remark}

\subsection{Non-circularity and non-overbreadth}\label{subsec:non-circularity-non-overbreadth}

Two proposition schemas record the definition-theoretic facts used by
the later theorem chain: the domain does not presuppose the answer to
the exact-six question, and it is structurally restrictive enough
for that question to remain substantive.

\begin{proposition}[Non-circularity schema for \FATCD{}]\label{prop:non-circularity}
Inspection of Definitions~\ref{def:raw-record}--\ref{def:fatcd} shows that
membership in \FATCD{} is not defined by requiring a decomposition/exhaustion record or
by requiring any derived role projection to exist. In particular:
\begin{enumerate}[label=(\roman*),topsep=2pt,itemsep=1pt]
  \item no clause of Definitions~\ref{def:raw-record},
    \ref{def:annotations}, \ref{def:finite}, \ref{def:closed}, or
    \ref{def:audited} quantifies over the role index \(i\) or requires a
    role-channel assignment;
  \item no membership clause excludes a candidate seventh-role feature by
    definition;
  \item the residual classification scheme (Definition~\ref{def:residual-scheme})
    is declarative and carries the categories (x) and (xi) as explicit
    failure options.
\end{enumerate}
\end{proposition}

\begin{proposition}[Non-overbreadth schema for \FATCD{}]\label{prop:non-overbreadth}
Inspection of Definitions~\ref{def:finite}--\ref{def:threshold-attachment}
shows that \FATCD{} is structurally restricted and does not admit arbitrary
syntactic material:
\begin{enumerate}[label=(\roman*),topsep=2pt,itemsep=1pt]
  \item \(\Raw{D}\) contains finitely many records and annotations, each
    record or annotation carries finite typed fields, and no infinite
    semantic, probabilistic, topological, geometric, causal, optimization, or
    resource structure is admitted unless finitely represented and bookkept as
    a raw feature, exactly as required by Definition~\ref{def:finite};
  \item every reference is bound to a declared record
    (Definition~\ref{def:closed});
  \item every claim is bookkept, and every promotion, translation, lift,
    forgetting, visibility, and empirical-bridge assertion carries the
    structured audit record required by
    Definition~\ref{def:audited};
  \item a role projection cannot be asserted without the full attachment
    data of Definition~\ref{def:threshold-attachment}.
\end{enumerate}
Consequently, arbitrary unaudited syntactic material cannot enter \FATCD{},
and mentioning a role name without the supporting threshold, witness, level,
collapse, and bookkeeping data does not by itself produce a derived role
projection.
\end{proposition}

\subsection{The scoped exact-six question over \FATCD{}}\label{subsec:scoped-exact-six-question}

We can now state the scoped exact-six question precisely.

\begin{definition}[Scoped exact-six question]\label{def:scoped-exact-six-question}
The \emph{scoped exact-six question} over \FATCD{} asks whether, for every
\(D \in \FATCD{}\), there exists at least one decomposition/exhaustion record
\(\Dec{D}\) whose role-channel structure consists of exactly the six typed
closure-role channels \Pone{}, \Ptwo{}, \Pthree{}, \Pfour{}, \Pfive{},
\Psix{}, each assigned exactly one of the five statuses later formalized in
Section~\ref{sec:decomposition-exhaustion}, with all residual features of
\(D\) classified under the scheme of Definition~\ref{def:residual-scheme}
without recourse to categories (x) or (xi). No uniqueness is built into this
question. The affirmative answer is denoted \(\ScopedExactSix{\FATCD{}}\).
\end{definition}

\begin{remark}[Channels, not active projections]\label{rem:channel-vs-activation-reminder}
An affirmative answer to the scoped exact-six question asserts the
existence of six role channels per description, not six active derived
projections per description. The distinction introduced in
Remark~\ref{rem:projection-vs-membership} and formalized in
Section~\ref{sec:decomposition-exhaustion} is preserved throughout.
\end{remark}

\begin{remark}[Relation to the Foundations theorem]\label{rem:fatcd-foundations-relation}
Proving \(\ScopedExactSix{\FATCD{}}\) is not a re-derivation of the
canonical unavoidability result of Foundations~I~\citep{Tsiokos2026SixBirdsFoundations},
and \(\ScopedExactSix{\FATCD{}}\) does not follow from that result
alone. The Foundations theorem derives the \Pone{}--\Psix{} vocabulary
under minimal hypotheses of composability with limited interface
access; the scoped exact-six claim asserts that, within the finite
audited typed closure regime of \FATCD{}, every admissible description
admits a six-channel decomposition/exhaustion record without a seventh
irreducible typed role in the covered scope. The two results share
vocabulary but address different questions.
\end{remark}
